\documentclass[../../main.tex]{subfiles}
\begin{document}

\section{群的自由积}\label{section:ch11-群的自由积}

在本节中，我们考虑那些未必满足交换性的群 $G$。此时我们将群 $G$ 的运算视为乘法，并且用 $1$ 表示 $G$ 的单位元，用 $x^{-1}$ 表示 $x$ 的逆元，用 $x^n$ 表示 $x$ 的 $n$-重乘积，用 $x^{-n}$ 表示 $x^{-1}$ 的 $n$-重乘积，用 $x^0$ 表示 $1$。

本节我们将讨论群的自由积这个概念，它对于任意群的作用相当于阿贝尔群的直和。

设 $G$ 是一个群。如果 $\{G_\alpha\}_{\alpha\in J}$ 是 $G$ 的一个子群族，如上节那样，我们称它们生成（generate）$G$，如果 $G$ 的每一个元素 $x$ 都能表示成群族 $G_\alpha$ 中所包含的元素的有限乘积。这意味着存在含于这些 $G_\alpha$ 的有限序列 $(x_1,\cdots,x_n)$ 使得 $x=x_1\cdots x_n$。这一序列称为群族 $G_\alpha$ 中（长度为 $n$）的字（word）。这时，我们说 $G$ 的元素 $x$ 由它表示（represent）。

注意，由于这里没有交换性的假定，在 $x$ 的表示中不能通过重组其因子而达到每一个因子组属于（不同的）$G_\alpha$ 的目的。然而，当 $x_i$ 与 $x_{i+1}$ 同属于一个 $G_\alpha$ 时，我们可以将其重组而得到一个长度为 $n-1$ 的字
\begin{align*}
(x_1,\cdots,x_{i-1},x_i x_{i+1},x_{i+2},\cdots,x_n),
\end{align*}
这个字也表示 $x$。此外，若 $x_i$ 等于 $1$，可以将其从序列中去掉，而得到表示 $x$ 的较短的字。

连续地实施上述化简运算，可以得到一个表示 $x$ 的字 $(y_1,\cdots,y_m)$ 使得 $y_i$ 和 $y_{i+1}$ 不包含于同一群 $G_\alpha$ 之中，并且对于每一个 $i$，$y_i\neq 1$。这样的字称为约化字（reduced word）。然而，若 $x$ 为 $G$ 的单位元，上述化简运算不再可行。这是由于此时 $x$ 可由字 $(a,a^{-1})$ 表示，这个字已简化到长度为 $1$ 的字 $(aa^{-1})$，从而同时消失！由此，为方便起见，我们约定：空集是表示 $G$ 中单位元（长度为 $0$）的约化字。据此，我们有：若群族 $G_\alpha$ 生成 $G$，则 $G$ 的每一个元素都可以用含于群族 $G_\alpha$ 的元素组成的约化字表示。

注意，若 $(x_1,\cdots,x_n)$ 和 $(y_1,\cdots,y_m)$ 分别为表示 $x$ 和 $y$ 的字，则 $(x_1,\cdots,x_n,y_1,\cdots,y_m)$ 为表示 $xy$ 的字。即使三个字中的前两个都是约化字，第三个也未必是约化字，除非 $x_n$ 与 $y_1$ 位于不同的 $G_\alpha$ 中。

\begin{definition}[群的自由积]\label{definition:ch11-free-product}
设 $G$ 是一个群，$\{G_\alpha\}_{\alpha\in J}$ 为生成 $G$ 的 $G$ 的一个子群族。假设当 $\alpha\neq\beta$ 时 $G_\alpha\cap G_\beta$ 仅含有单位元。称 $G$ 为群族 $G_\alpha$ 的自由积（free product），如果对于任何一个 $x\in G$，都唯一地存在含于群族 $G_\alpha$ 并且表示 $x$ 的字。此时，记作
\begin{align*}
G=\prod_{\alpha\in J}^{*}G_\alpha,
\end{align*}
对于有限的情形，也记作 $G=G_1*\cdots*G_n$。
\end{definition}

设 $G$ 是群族 $G_\alpha$ 的自由积，$(x_1,\cdots,x_n)$ 为含于群 $G_\alpha$ 中的元素构成的字，对于每一个 $i$ 有 $x_i\neq 1$。则对于每一个 $i$，存在一个唯一 $\alpha_i$ 使得 $x_i\in G_{\alpha_i}$。我们可以将“它是一个约化字”简单地说成：对于每一个 $i$ 有 $\alpha_i\neq\alpha_{i+1}$。

设群族 $G_\alpha$ 生成 $G$，其中当 $\alpha\neq\beta$ 时有 $G_\alpha\cap G_\beta=\{1\}$。为了使得 $G$ 为这些群的自由积，只需判定由空字给出的 $1$ 的表示是唯一的。其原因如下：假设这个较弱的条件成立，并且假设 $(x_1,\cdots,x_n)$ 和 $(y_1,\cdots,y_m)$ 是表示 $G$ 中同一元素 $x$ 的两个约化字。设 $\alpha_i$ 和 $\beta_i$ 是使得 $x_i\in G_{\alpha_i}$ 及 $y_i\in G_{\beta_i}$ 的指标。由于
\begin{align*}
x_1\cdots x_n=x=y_1\cdots y_m,
\end{align*}
从而有
\begin{align*}
(y_m^{-1},\cdots,y_1^{-1},x_1,\cdots,x_n)
\end{align*}
表示 $1$。这样的字必定是可约化的，由此可见必有 $\alpha_1=\beta_1$，从而这个字可约化为
\begin{align*}
(y_m^{-1},\cdots,y_1^{-1}x_1,\cdots,x_n).
\end{align*}
重复以上讨论可见，这个字也一定是可约化的，从而有 $y_1^{-1}x_1=1$。那么 $x_1=y_1$，由此，$1$ 可以用
\begin{align*}
(y_m^{-1},\cdots,y_2^{-1},x_2,\cdots,x_n)
\end{align*}
表示。不断重复上述讨论可以得出 $m=n$ 并且对于每一个 $i$ 有 $x_i=y_i$。

\begin{example}\label{example:ch11-permutation-group}
考虑由集合 $\{0,1,2\}$ 到自身的一一映射所构成的群 $P$。对于 $i=1,2$，定义 $P$ 中的元素 $\pi_i$ 满足条件：$\pi_i(i)=i-1$，$\pi_i(i-1)=i$ 以及当 $j\neq 1,2$ 时，$\pi_i(j)=j$。则 $\pi_i$ 生成 $P$ 的一个子群 $G_i$，其阶为 $2$。可以验证：$G_1$ 和 $G_2$ 生成 $P$。但 $P$ 不是它们的自由积。比如，两个约化字 $(\pi_1,\pi_2,\pi_1)$ 和 $(\pi_2,\pi_1,\pi_2)$ 表示 $P$ 中的同一元素。
\end{example}

自由积也满足类似于直和的扩展条件（extension condition）。

\begin{lemma}[自由积的扩展条件]\label{lemma:ch11-extension-condition}
设 $G$ 是一个群，$\{G_\alpha\}$ 为 $G$ 的子群族。若 $G$ 为群族 $G_\alpha$ 的自由积，则 $G$ 满足以下条件：

$(*)$ 对于任何群 $H$ 以及任何同态族 $h_\alpha:G_\alpha\to H$，存在一个同态 $h:G\to H$，使得对于每一个 $\alpha$，同态 $h$ 在 $G_\alpha$ 上的限制等于 $h_\alpha$。

此外，$h$ 是唯一的。
\end{lemma}

该引理的逆也是成立的，只是其证明不像直和的情况那样简单。我们稍后再来讨论。

\begin{proof}
给定 $x\in G$，$x\neq 1$。设 $(x_1,\cdots,x_n)$ 是表示 $x$ 的约化字。若 $h$ 存在，它必定满足等式
\begin{align}
h(x)=h(x_1)\cdots h(x_n)=h_{\alpha_1}(x_1)\cdots h_{\alpha_n}(x_n),\label{align:ch11-free-product-hom}
\end{align}
其中 $\alpha_i$ 为满足条件 $x_i\in G_{\alpha_i}$ 的指标。由此 $h$ 是唯一的。

为了证明 $h$ 的存在性，我们定义 $h$ 满足条件：当 $x\neq 1$ 时，$h$ 由 \eqref{align:ch11-free-product-hom} 定义，并且令 $h(1)=1$。由于由约化字给出的 $x$ 的表示是唯一的，所以 $h$ 的定义确切。我们证明 $h$ 是一个同态。

首先证明一个预备性结论。任意给定一个具有正长度的由含于群族 $G_\alpha$ 的元素表示的字 $w=(x_1,\cdots,x_n)$，定义 $\phi(w)$ 为 $H$ 中满足等式
\begin{align}
\phi(w)=h_{\alpha_1}(x_1)\cdots h_{\alpha_n}(x_n)\label{align:ch11-free-product-phi}
\end{align}
的元素，其中 $\alpha_i$ 是使得 $x_i\in G_{\alpha_i}$ 的任何一个指标。当 $x_i=1$ 时 $\alpha_i$ 是唯一的，因此 $\phi$ 的定义确切。若 $w$ 是空字，令 $\phi(w)$ 表示 $H$ 上的单位元。我们证明：若 $w'$ 是 $w$ 经某一化简运算后所得到的字，则 $\phi(w')=\phi(w)$。

首先假设 $w'$ 是从字 $w$ 中删除 $x_i=1$ 后得到的字，那么由 $h_{\alpha_i}(x_i)=1$ 推得 $\phi(w')=\phi(w)$。

其次，假设 $\alpha_i=\alpha_{i+1}$ 且
\begin{align*}
w'=(x_1,\cdots,x_i x_{i+1},\cdots,x_n).
\end{align*}
事实上，等式
\begin{align*}
h_\alpha(x_i)h_\alpha(x_{i+1})=h_\alpha(x_i x_{i+1})
\end{align*}
蕴涵着 $\phi(w)=\phi(w')$，其中 $\alpha=\alpha_i=\alpha_{i+1}$。

若 $w$ 为群族 $G_\alpha$ 中表示 $x$ 的任何一个字，则容易推得 $h(x)=\phi(w)$。因为由 $h$ 的定义，等式对任何约化字 $w$ 成立，并且化简运算不改变 $\phi$ 的值。

现在来证明 $h$ 是一个同态。设 $w=(x_1,\cdots,x_n)$ 及 $w'=(y_1,\cdots,y_m)$ 分别是表示 $x$ 和 $y$ 的字。设 $(w,w')$ 为表示 $xy$ 的字 $(x_1,\cdots,x_n,y_1,\cdots,y_m)$。由等式 \eqref{align:ch11-free-product-phi} 得到 $\phi(w,w')=\phi(w)\phi(w')$，从而 $h(xy)=h(x)h(y)$。
\end{proof}

现在我们考虑以下问题：对于任何群的一个族 $\{G_\alpha\}$，给出一个群 $G$，它包含所有子群 $G_\alpha'$，这个子群同构于 $G_\alpha$，使得 $G$ 是子群族 $G_\alpha'$ 的自由积。我们将看到这是可以做到的，但需要引入一个外自由积的概念。

\begin{definition}[外自由积]\label{definition:ch11-external-free-product}
设 $\{G_\alpha\}_{\alpha\in J}$ 为群的一个加标族。设 $G$ 是一个群，$i_\alpha:G_\alpha\to G$ 为单同态的一个族，使得 $G$ 为群族 $i_\alpha(G_\alpha)$ 的自由积。则称 $G$ 为群族 $G_\alpha$ 相对于单同态族 $i_\alpha$ 的外自由积（external free product）。
\end{definition}

当然，群 $G$ 不是唯一的，然而我们将证明在不区别同构的群的意义下它是唯一的。$G$ 的构造比外直和的构造要困难一些。

\begin{theorem}[外自由积的存在性]\label{theorem:ch11-external-free-product-existence}
对于给定的群的一个族 $\{G_\alpha\}_{\alpha\in J}$，存在一个群 $G$ 及单同态的一个族 $i_\alpha:G_\alpha\to G$，使得 $G$ 是群族 $i_\alpha(G_\alpha)$ 的自由积。
\end{theorem}

\begin{proof}
为方便起见，我们假定群族 $G_\alpha$ 作为集合族是两两无交的。（这容易办到，因为必要的话，可以用 $G_\alpha\times\{\alpha\}$ 替换 $G_\alpha$。）

如前所述，我们将包含于群族 $G_\alpha$ 中的元素定义一个（长度为 $n$ 的）字（word），即 $\bigcup G_\alpha$ 中元素的 $n$-元素串 $w=(x_1,\cdots,x_n)$。这个字称为约化字，如果对于每一个 $i$，$\alpha_i\neq\alpha_{i+1}$，其中 $\alpha_i$ 为使得 $x_i\in G_{\alpha_i}$ 的指标，并且对于每一个 $i$，$x_i$ 不是 $G_{\alpha_i}$ 的单位元。我们定义空集是唯一的具有 $0$ 长度的约化字。注意：我们还没有给出含所有子群 $G_\alpha$ 的群 $G$，因此目前还不能谈到“表示”$G$ 的元素的字。

设 $W$ 表示群族 $G_\alpha$ 中所有约化字所组成的集合。令 $P(W)$ 表示所有一一映射 $\pi:W\to W$ 的集合，则 $P(W)$ 是一个群，它以映射的复合运算为群的运算。$P(W)$ 的某一个子群 $G$ 将是我们所需要的群。

\textbf{第一步。} 对每一指标 $\alpha$ 及每一个 $x\in G_\alpha$，我们定义一个满足以下条件的映射 $\pi_x:W\to W$：

(1) 以 $1_\alpha$ 表示 $G_\alpha$ 的单位元。若 $x=1_\alpha$，则 $\pi_x$ 为 $W$ 中的恒等映射。

(2) 若 $x,y\in G_\alpha$ 并且 $z=xy$，则 $\pi_z=\pi_x\circ\pi_y$。

具体定义如下：设 $x\in G_\alpha$。先明确几个记号，$w=(x_1,\cdots,x_n)$ 总是表示 $W$ 中的非空元素，$\alpha_1$ 表示满足 $x_1\in G_{\alpha_1}$ 的指标。若 $x\neq 1_\alpha$，定义 $\pi_x$ 如下：

(i) $\pi_x(\varnothing)=(x)$，

(ii) $\pi_x(w)=(x,x_1,\cdots,x_n)$ 若 $\alpha_1\neq\alpha$，

(iii) $\pi_x(w)=(xx_1,\cdots,x_n)$ 若 $\alpha_1=\alpha$ 且 $x_1\neq x^{-1}$，

(iv) $\pi_x(w)=(x_2,\cdots,x_n)$ 若 $\alpha_1=\alpha$ 且 $x_1=x^{-1}$。

如果 $x=1_\alpha$，定义 $\pi_x$ 为 $W$ 的恒等映射。

注意对每一种情形，$\pi_x$ 的值都是一个约化字，即是 $W$ 中的元素。对于情形 (i) 和 (ii)，在 $\pi_x$ 的作用下字的长度增加；对于情形 (iii)，在 $\pi_x$ 的作用下保持字的长度不变；对于情形 (iv)，在 $\pi_x$ 的作用下字的长度减少。对于情形 (iv)，$\pi_x$ 作用在长度为 $1$ 的字 $w$ 上，其结果为一个空字。

\textbf{第二步。} 我们证明：如果 $x,y\in G_\alpha$ 并且 $z=xy$，则 $\pi_z=\pi_x\circ\pi_y$。

如果 $x$ 或者 $y$ 等于 $1_\alpha$，结论显然成立，因为此时 $\pi_x$ 或者 $\pi_y$ 为恒等映射。以下假设 $x\neq 1_\alpha$ 和 $y\neq 1_\alpha$。我们来计算 $\pi_z$ 以及 $\pi_x\circ\pi_y$ 在约化字 $w$ 处的值。需要考虑四种情形。

(i) 设 $w$ 是一个空字。我们有 $\pi_y(\varnothing)=(y)$。如果 $z=1_\alpha$，则 $y=x^{-1}$ 并且根据 (iv) 有 $\pi_x\pi_y(\varnothing)=\varnothing$，同时 $\pi_z(\varnothing)$ 也是空集，因为 $\pi_z$ 为恒等映射。如果 $z\neq 1_\alpha$，则
\begin{align*}
\pi_x\pi_y(\varnothing)=(xy)=(z)=\pi_z(\varnothing).
\end{align*}

对于其他情形，假设 $w=(x_1,\cdots,x_n)$，其中 $x_1\in G_{\alpha_1}$。

(ii) 假设 $\alpha\neq\alpha_1$。则 $\pi_y(w)=(y,x_1,\cdots,x_n)$。若 $z=1_\alpha$，则 $y=x^{-1}$ 且根据 (iv) 有 $\pi_x\pi_y(w)=(x_1,\cdots,x_n)$，同时由于 $\pi_z$ 为恒等映射，$\pi_z(w)=w$。如果 $z\neq 1_\alpha$，则
\begin{align*}
\pi_x\pi_y(w)=(xy,x_1,\cdots,x_n)=(z,x_1,\cdots,x_n)=\pi_z(w).
\end{align*}

(iii) 设 $\alpha=\alpha_1$ 和 $yx_1\neq 1_\alpha$。则 $\pi_y(w)=(yx_1,x_2,\cdots,x_n)$。如果 $xyx_1=1_\alpha$，则 $\pi_x\pi_y(w)=(x_2,\cdots,x_n)$，同时，由于 $zx_1=xyx_1=1_\alpha$，$\pi_z(w)$ 等于 $w$。如果 $xyx_1\neq 1_\alpha$，则
\begin{align*}
\pi_x\pi_y(w)=(xyx_1,x_2,\cdots,x_n)=(zx_1,x_2,\cdots,x_n)=\pi_z(w).
\end{align*}

(iv) 最后，假设 $\alpha=\alpha_1$ 和 $yx_1=1_\alpha$。则 $\pi_y(w)=(x_2,\cdots,x_n)$，当 $n=1$ 时这是一个空字。这时，
\begin{align*}
\pi_x\pi_y(w)=(x,x_2,\cdots,x_n)=(x(yx_1),x_2,\cdots,x_n)=(zx_1,x_2,\cdots,x_n)=\pi_z(w).
\end{align*}

\textbf{第三步。} 映射 $\pi_x$ 是 $P(W)$ 中的一个元素，并且由 $i_\alpha(x)=\pi_x$ 定义的映射 $i_\alpha:G_\alpha\to P(W)$ 是一个单同态。

为了证明 $\pi_x$ 是一个一一映射，注意到如果 $y=x^{-1}$，则条件 (1) 和 (2) 蕴涵着 $\pi_y\circ\pi_x$ 与 $\pi_x\circ\pi_y$ 都是 $W$ 上的恒等映射。由此，$\pi_x$ 属于 $P(W)$。事实上，条件 (2) 蕴涵着 $i_\alpha$ 是一个同态。为了证明 $i_\alpha$ 是一个单同态，只要注意如果 $x\neq 1_\alpha$，则 $\pi_x(\varnothing)=(x)$，从而 $\pi_x$ 不是 $W$ 上的恒等映射。

\textbf{第四步。} 设 $G$ 是由群族 $G_\alpha'=i_\alpha(G_\alpha)$ 生成的 $P(W)$ 的一个子群。我们证明 $G$ 是群族 $G_\alpha'$ 的自由积。

首先，我们证明若 $\alpha\neq\beta$，则 $G_\alpha'\cap G_\beta'$ 仅含有单位元。设 $x\in G_\alpha$ 且 $y\in G_\beta$。假定 $\pi_x$ 和 $\pi_y$ 都不是 $W$ 上的恒等映射，并且证明 $\pi_x\neq\pi_y$。这是容易的，因为 $\pi_x(\varnothing)=(x)$ 而 $\pi_y(\varnothing)=(y)$，从而它们是不同的字。

其次，我们证明在群族 $G_\alpha'$ 中没有非空的约化字
\begin{align*}
w'=(\pi_{x_1},\cdots,\pi_{x_n})
\end{align*}
表示 $G$ 中的单位元。设 $\alpha_i$ 是使得 $x_i\in G_{\alpha_i}$ 的指标，则 $\alpha_i\neq\alpha_{i+1}$ 并且对于所有 $i$，$x_i\neq 1_{\alpha_i}$。这时，
\begin{align*}
\pi_{x_1}(\pi_{x_2}(\cdots(\pi_{x_n}(\varnothing))))=(x_1,\cdots,x_n),
\end{align*}
从而，由 $w'$ 表示的 $G$ 中的元素不是 $P(W)$ 的单位元。
\end{proof}

尽管上述关于自由积的存在性的证明的正确性是无可非议的，然而它却有一个缺陷，因为它没有给我们提供一个简便的办法来考虑自由积中的元素。在许多场合，这个缺陷对我们无所谓，因为许多应用的关键在于它的存在性。然而，人们希望对自由积的具体模型有一个更为清晰的了解。

就外直和而言，我们做到了这一点。阿贝尔群族 $G_\alpha$ 的外直和由笛卡儿积 $\prod G_\alpha$ 中的这样一些元素 $(x_\alpha)$ 所构成：仅有有限多个 $x_\alpha\neq 0_\alpha$。每一个群 $G_\beta$ 同构于由这样一些元素 $(x_\alpha)$ 所构成的子群 $G_\beta'$，这些元素对于所有的 $\alpha\neq\beta$ 有 $x_\alpha=0_\alpha$。

对于自由积是否也有这样一个简单的模型呢？回答是肯定的。在上述证明的最后一步，我们证明了：若 $(\pi_{x_1},\cdots,\pi_{x_n})$ 是 $G_\alpha'$ 中的一个约化字，则
\begin{align*}
\pi_{x_1}(\pi_{x_2}(\cdots(\pi_{x_n}(\varnothing))))=(x_1,\cdots,x_n).
\end{align*}
此等式蕴涵着如果 $\pi$ 是 $P(W)$ 中属于自由积 $G$ 的任何一个元素，则对应法则 $\pi\mapsto\pi(\varnothing)$ 定义了一个 $G$ 与集合 $W$ 之间的一一对应！进而，如果 $\pi$ 和 $\pi'$ 为 $G$ 中的两个元素，使得
\begin{align*}
\pi(\varnothing)=(x_1,\cdots,x_n)\quad\text{和}\quad\pi'(\varnothing)=(y_1,\cdots,y_k),
\end{align*}
则 $\pi(\pi'(\varnothing))$ 是一个将字 $(x_1,\cdots,x_n,y_1,\cdots,y_k)$ 不断约化而得到的字！

这就为我们提供了一种了解 $G$ 的方法。可以将 $G$ 简单地视为集合 $W$，其上两个字的乘法运算就是将其连接并化简结果。单位元对应于空字。并且每一个群 $G_\beta$ 是 $W$ 的一个子集，它由空集以及所有形如 $(x)$ 的长度为 $1$ 的字所组成，其中 $x\in G_\beta$，$x\neq 1_\beta$。

一个问题马上来了：为什么不按照如上的说法去定义自由积呢？这比考虑 $W$ 的置换群 $P(W)$ 似乎是一种更为简洁的方式。对上述问题的回答是：如果以此种说法作为自由积的定义，则群公理的验证将变得十分困难，特别是结合律的验证将异常复杂。相比较而言，前面给出的关于自由积的存在性的证明更为简洁和优美。

对关于通常的自由积的扩展条件稍加修改，便可得到关于外自由积的扩展条件。

\begin{theorem}[外自由积的扩展条件]\label{theorem:ch11-external-extension-condition}
设 $\{G_\alpha\}$ 是群的一个族，$G$ 是一个群，$i_\alpha:G_\alpha\to G$ 为同态的一个族。如果每一个 $i_\alpha$ 都是一个单同态，并且 $G$ 为群族 $i_\alpha(G_\alpha)$ 的自由积，则 $G$ 满足以下条件：

对于任何一个群 $H$，以及任何同态的一个族 $h_\alpha:G_\alpha\to H$，存在一个同态 $h:G\to H$，使得对于每一个 $\alpha$ 都有 $h\circ i_\alpha=h_\alpha$。

此外，$h$ 是唯一的。
\end{theorem}

由此得到的一个结论是自由积的唯一性定理，它的证明与直和情形的证明非常相似，我们将其留给读者。

\begin{theorem}[自由积的唯一性]\label{theorem:ch11-uniqueness-of-free-product}
设 $\{G_\alpha\}_{\alpha\in J}$ 是群的一个族。假设 $G$ 和 $G'$ 都是群，$i_\alpha:G_\alpha\to G$ 和 $i_\alpha':G_\alpha\to G'$ 都是单同态族，使得 $\{i_\alpha(G_\alpha)\}$ 和 $\{i_\alpha'(G_\alpha)\}$ 分别生成 $G$ 和 $G'$。如果 $G$ 和 $G'$ 分别具有前述引理中陈述的扩展性质，则有一个唯一的同构 $\phi:G\to G'$，使得对于每一个 $\alpha$ 有 $\phi\circ i_\alpha=i_\alpha'$。
\end{theorem}

最后，我们能够证明扩展条件决定了自由积，即证明引理 68.1 和 68.3 的逆。

\begin{lemma}\label{lemma:ch11-converse-extension}
设 $\{G_\alpha\}_{\alpha\in J}$ 是群的一个族。假设 $G$ 是一个群，$i_\alpha:G_\alpha\to G$ 是同态的一个族。如果定理 68.3 中的扩展条件成立，则每一个 $i_\alpha$ 都是一个单同态，并且 $G$ 是群族 $i_\alpha(G_\alpha)$ 的自由积。
\end{lemma}

\begin{proof}
我们首先证明每一个 $i_\alpha$ 是一个单同态。任意给定一个指标 $\beta$，记 $H=G_\beta$。设 $h_\alpha:G_\alpha\to H$ 满足条件：当 $\alpha=\beta$ 时为恒等映射，当 $\alpha\neq\beta$ 时为平凡同态。设 $h:G\to H$ 为由扩展条件给出的同态。则 $h\circ i_\beta=h_\beta$，从而 $i_\beta$ 为单射。

根据定理 68.2，存在一个群 $G'$ 和单同态的一个族 $i_\alpha':G_\alpha\to G'$ 使得 $G'$ 是群 $i_\alpha'(G_\alpha)$ 的自由积。$G$ 和 $G'$ 都具有引理 68.3 中的扩展性质。则前面的定理蕴涵着：存在一个同构 $\phi:G\to G'$ 使得 $\phi\circ i_\alpha=i_\alpha'$。由此立即得到 $G$ 为群族 $i_\alpha(G_\alpha)$ 的自由积。
\end{proof}

我们再来证明与推论 67.2 和 67.3 类似的两个结果。

\begin{corollary}\label{corollary:ch11-free-product-associativity}
设 $G=G_1*G_2$，其中 $G_1$ 和 $G_2$ 分别是子群族 $\{H_\alpha\}_{\alpha\in J}$ 和 $\{H_\beta\}_{\beta\in K}$ 的自由积。如果指标集 $J$ 与 $K$ 无交，则 $G$ 为子群族 $\{H_\gamma\}_{\gamma\in J\cup K}$ 的自由积。
\end{corollary}

\begin{proof}
这个证明几乎是推论 67.2 的证明的翻版。
\end{proof}

这个结论特别蕴涵着
\begin{align*}
G_1*G_2*G_3=G_1*(G_2*G_3)=(G_1*G_2)*G_3.
\end{align*}

为了表述下一个定理，我们需回忆一下群论中的某些概念。设 $x$ 和 $y$ 都是群 $G$ 中的元素，称 $y$ 共轭（conjugate）于 $x$，如果对于每一个 $c\in G$，$y=cxc^{-1}$ 成立。群 $G$ 的正规子群是一个含有群中所有元素的所有共轭元的子群。

设 $S$ 为 $G$ 的一个子集，我们考虑 $G$ 中包含 $S$ 的所有正规子群的交 $N$。易见，$N$ 自身也是一个 $G$ 的正规子群，称为 $G$ 中包含 $S$ 的最小的正规子群（least normal subgroup）。

\begin{theorem}\label{theorem:ch11-quotient-of-free-product}
设 $G=G_1*G_2$。设 $N_i$ 为 $G_i$ 的正规子群，$i=1,2$。如果 $N$ 是 $G$ 中包含 $N_1$ 和 $N_2$ 的最小正规子群，则
\begin{align*}
G/N\cong(G_1/N_1)*(G_2/N_2).
\end{align*}
\end{theorem}

\begin{proof}
内射与投影同态的复合
\begin{align*}
G_1\longrightarrow G_1*G_2\longrightarrow(G_1*G_2)/N
\end{align*}
将 $N_1$ 映为单位元，从而它诱导了一个同态
\begin{align*}
i_1:G_1/N_1\longrightarrow(G_1*G_2)/N.
\end{align*}
类似地，内射与投影同态的复合诱导了一个同态
\begin{align*}
i_2:G_2/N_2\longrightarrow(G_1*G_2)/N.
\end{align*}
我们来证明 $i_1$ 和 $i_2$ 满足引理 68.5 中的扩展条件。由此可见，$i_1$ 和 $i_2$ 是单同态并且 $(G_1*G_2)/N$ 是相对于这两个单同态的 $G_1/N_1$ 与 $G_2/N_2$ 的外自由积。

设 $h_1:G_1/N_1\to H$ 和 $h_2:G_2/N_2\to H$ 为任意两个同态，则关于 $G_1*G_2$ 的扩展条件蕴涵着存在一个从 $G_1*G_2$ 到 $H$ 的同态，它等于投射与 $G_i$ 上的映射 $h_i$ 的复合
\begin{align*}
G_i\longrightarrow G_i/N_i\longrightarrow H,
\end{align*}
其中 $i=1,2$。此同态将 $N_1$ 与 $N_2$ 中的元素映射为单位元，因此，它的核包含 $N$。从而，它诱导了一个同态 $h:(G_1*G_2)/N\to H$，满足 $h_1=h\circ i_1$ 及 $h_2=h\circ i_2$。
\end{proof}

\begin{corollary}\label{corollary:ch11-quotient-by-normal}
如果 $N$ 是 $G_1*G_2$ 中包含 $G_1$ 的最小正规子群，则 $(G_1*G_2)/N\cong G_2$。
\end{corollary}

在进一步的讨论中，我们将经常用到“最小正规子群”的概念。显然，如果 $N$ 是 $G$ 中包含子集 $S$ 的最小正规子群，则 $N$ 包含 $S$ 以及 $S$ 的所有元素的所有共轭元。为以后使用方便，我们来验证这些元素实际上生成 $N$。

\begin{lemma}\label{lemma:ch11-least-normal-subgroup}
设 $S$ 为群 $G$ 的一个子集。如果 $N$ 为 $G$ 中包含 $S$ 的最小正规子群，则 $S$ 中所有元素的所有共轭元生成 $N$。
\end{lemma}

\begin{proof}
设 $N'$ 是由所有 $S$ 的元素的所有共轭元生成的 $G$ 的子群。我们有 $N'\subset N$。为验证相反的包含关系，只需证明 $N'$ 是 $G$ 的正规子群。任意选取 $x\in N'$ 以及 $c\in G$，我们证明 $cxc^{-1}\in N'$。

可以将 $x$ 写成 $x=x_1x_2\cdots x_n$ 的形式，其中每一个 $x_i$ 都共轭于 $S$ 中的某一个元素 $s_i$。这时 $cxc^{-1}$ 也共轭于 $s_i$。因为
\begin{align*}
cxc^{-1}=(cx_1c^{-1})(cx_2c^{-1})\cdots(cx_nc^{-1}),
\end{align*}
从而 $cxc^{-1}$ 是 $S$ 的元素的共轭元的乘积，因此 $cxc^{-1}\in N'$。
\end{proof}

\subsection*{习题}

\begin{exercise}\label{exercise:ch11-68-01}
验证例 1 中的细节。
\end{exercise}

\begin{exercise}\label{exercise:ch11-68-02}
设 $G=G_1*G_2$，其中 $G_1$ 和 $G_2$ 都是非平凡群。
\begin{enumerate}[(a)]
\item\label{item:ch11-68-02-a} 证明 $G$ 不是阿贝尔群。
\item 对于 $x\in G$，定义 $x$ 的长度为 $G_1$ 和 $G_2$ 中表示 $x$ 的元素的唯一的约化字的长度。证明：如果 $x$ 的长度为偶数（不小于 $2$），则 $x$ 不是无限阶的。如果 $x$ 的长度为奇数，则 $x$ 共轭于一个具有较短长度的元素。
\item 证明 $G$ 中的有限阶元素是 $G_1$ 和 $G_2$ 中的有限阶元素及其共轭元。
\end{enumerate}
\end{exercise}

\begin{exercise}\label{exercise:ch11-68-03}
设 $G=G_1*G_2$。对于 $c\in G$，令 $cG_1c^{-1}$ 表示所有形如 $cxc^{-1}$ 的元素，其中 $x\in G_1$。它是 $G$ 的一个子群，证明它与 $G_2$ 的交只含有单位元。
\end{exercise}

\begin{exercise}\label{exercise:ch11-68-04}
证明定理 68.4。
\end{exercise}

\end{document}